% ======================================================================
\section{2.4.4 Dowód Twierdzenia 5}
% ======================================================================

\begin{frame}{Algorytm}
  \begin{enumerate}
    \item Policz kodowania wszystkich pasów wierszy i pasów kolumn (2.4.1).
    \item<2-> Każda chciana pozycja $(I,J)\in W$ leży w \textbf{jednym} kafelku $T$ i jest wskazana przez \textbf{jeden} napis
      wyjściowy kafelka $T$. Jego zbiór wewnętrzny to zbiór $Q$ iloczynu blokowego zawierającego $(I,J)$, a wiersz i kolumna to wiersz
      i kolumna $(I,J)$ wewnątrz tego iloczynu blokowego (2.3.3, 2.3.4).
    \item<3-> Dla kafelka $T$ niech $W_T$ będzie zbiorem napisów wyjściowych chcianych pozycji w $T$.
    \item<4-> Dla każdego kafelka $T$ z $W_T\ne\emptyset$ uruchom $\Pruned(W_T)$.
    \item<5-> Dla każdego $(I,J)\in W$ zwróć wartość przy jego napisie wyjściowym.
  \end{enumerate}
  \uncover<6->{\vspace{0.5em}\textbf{Poprawność:} z Lematu 9 (napis wskazuje właściwy wpis iloczynu blokowego, czyli $(XY)[I,J]$) i Lematu 10
    ($\Pruned$ zwraca $\Mult$ na $W_T$).}
  \uncover<7->{
  \begin{center}
    \begin{tikzpicture}[b/.style={draw,rounded corners,fill=blue!8,font=\scriptsize,align=center,minimum height=0.9cm},
        arr/.style={-{Stealth},thick}]
      \node[b] (e) {kodowania\\ pasów};
      \node[b,right=0.6cm of e] (w) {$W\to W_T$\\ (napisy)};
      \node[b,right=0.6cm of w,fill=red!12] (p) {$\Pruned(W_T)$\\ dla każdego $T$};
      \node[b,right=0.6cm of p] (o) {odczyt\\ $(XY)[I,J]$};
      \draw[arr] (e)--(w); \draw[arr] (w)--(p); \draw[arr] (p)--(o);
    \end{tikzpicture}
  \end{center}}
\end{frame}

\begin{frame}{Konsekwencja (a) założenia $N\ge D^{18}$}
  Ustalamy $D=4^m$ i $L=19m$. Mamy $D^{18}=4^{18m}$, więc potęgi $m$-te porównujemy, porównując ich podstawy.
  \pause
  \begin{block}{(a) Kafelkowanie: $K_0N_0\le N$, więc dopełnienie co najwyżej podwaja $N$}
    Z ogólnego oszacowania $\binom Lm\le(eL/m)^m$: $K=\binom{19m}{m}\le(19e)^m$. Zatem
    \[
      KN_0\le\big(19e\cdot3^{18}\big)^m\le\big(4^{18}\big)^m,\qquad\text{bo}\quad
      19e\cdot3^{18}\approx2{,}001\cdot10^{10}<4^{18}\approx6{,}872\cdot10^{10},
    \]
    a stąd $K_0N_0\le KN_0\le D^{18}\le N$. Dalej $N$ oznacza dopełniony rozmiar.
  \end{block}
\end{frame}

\begin{frame}{Konsekwencja (b) założenia $N\ge D^{18}$}
  \begin{block}{(b) Kodowania są tanie}
    Wykonujemy $2N/(K_0N_0)$ kodowań po $O(10^L)$ operacji. Żeby zmieścić się w budżecie $O(2^{-m/9}N^2)$, potrzeba
    \[
      10^L\le2^{-m/9}\,N\,N_0.\tag{6}
    \]
    Zachodzi:
    $\dfrac{10^L\,2^{m/9}}{N_0}=\Big(\dfrac{10^{19}\cdot2^{1/9}}{3^{18}}\Big)^m\le(4^{18})^m=D^{18}\le N$, \
    bo $\dfrac{10^{19}\cdot2^{1/9}}{3^{18}}\approx2{,}788\cdot10^{10}<4^{18}$.
  \end{block}
  \pause
  Wtedy łączny koszt kodowań to
  \[
    \frac{2N}{K_0N_0}\cdot O(10^L)\overset{(6)}{\le}\frac{2N}{K_0N_0}\cdot O\big(2^{-m/9}NN_0\big)=O\Big(\frac{2^{-m/9}N^2}{K_0}\Big)=O\big(2^{-m/9}N^2\big).
  \]
\end{frame}

\begin{frame}{Koszt}
  \begin{itemize}
    \item \textbf{Kodowania:} z (6)
      \[
        \frac{2N}{K_0N_0}\cdot O(10^L)\le\frac{2N}{K_0N_0}\cdot O\big(2^{-m/9}NN_0\big)=O\big(2^{-m/9}N^2\big).
      \]
    \item<2-> \textbf{Chciane wpisy:} każda pozycja z $W$ trafia do jednego $W_T$, więc
      $\sum_T|W_T|\le|W|\le N^2/\sqrt D=2^{-m}N^2$. Kafelków jest co najwyżej $4N^2/M$, więc $\sum_TM\le4N^2$.
    \item<3-> \textbf{Jedno wywołanie $\Pruned$ ze zbiorem $S$:} $O(L|S|)$ operacji. Każdy krok obsługuje każdy napis z $S$ stałą liczbę razy,
      a napis ma długość $L$. Zbiory trzymamy jako posortowane listy, więc slice to spójny fragment, a suma slice'ów to scalanie.
    \item<4-> Z Lematu 10 (łączny rozmiar zbiorów $\le(L+1)|\Leaves|$) wychodzi $O(L^2)=O(\log^2D)$ operacji \textbf{na liść} z $\Leaves(W_T)$.
  \end{itemize}
  \uncover<5->{
  Z Lematu 11 zsumowanego po kafelkach:
  \[
    O\Big(L^2\sum_T|\Leaves(W_T)|\Big)
    \le O\Big(L^2\,2^{-m/9}\Big(\underbrace{2^m|W|}_{\le N^2}+\underbrace{2\textstyle\sum_TM}_{\le8N^2}\Big)\Big)
    =O\big(L^2\,2^{-m/9}N^2\big),
  \]
  \alert{niezależnie od tego, jak $W$ rozkłada się na kafelki}.}
\end{frame}

\begin{frame}{Koszt: drobiazgi i wynik końcowy}
  \begin{itemize}
    \item \textbf{Drobiazgi:} wypisanie $K_0^2$ zbiorów $Q$, znalezienie napisów wyjściowych i sortowanie zbiorów $W_T$ kosztują $O(L)$
      na zbiór $Q$, na kafelek i na pozycję z $W$. Razem to $O(L\,2^{-m/9}N^2)$, bo:
      \begin{itemize}
        \item $K\le N$,
        \item kafelków jest co najwyżej $4N^2/N_0$,
        \item $|W|\le2^{-m}N^2$,
        \item $N,N_0\ge2^{m/9}$.
      \end{itemize}
  \end{itemize}
  \pause
  \begin{block}{Razem}
    Ponieważ $L=19\log_4D$ i $2^{-m/9}=D^{-1/18}$:
    \[
      O\big(2^{-m/9}N^2\big)+O\big(L^2\,2^{-m/9}N^2\big)+O\big(L\,2^{-m/9}N^2\big)
      =O\!\left(\frac{N^2\log^2D}{D^{1/18}}\right).\qquad\blacksquare
    \]
  \end{block}
\end{frame}

\begin{frame}{Rozmiar słowa}
  Każda liczba w algorytmie jest liczbą całkowitą o module $N^{O(1)}$:
  \begin{itemize}
    \item Wpisy wejścia mają taki rozmiar z założenia.
    \item<2-> Każda liczba w kodowaniu (lub policzona po drodze do niego) to kombinacja $\pm1$ co najwyżej $7^L$ wpisów wejścia,
      bo krok (2) tylko dodaje i odejmuje slice'y.
    \item<3-> Każda wartość policzona przez $\Pruned$ to suma co najwyżej $10^L$ iloczynów dwóch zakodowanych liczb (Lemat 7).
    \item<4-> Z (6): $7^L\le10^L\le2^{-m/9}NN_0\le N^2$.
    \item<5-> Indeksy i pozycje na listach nie przekraczają liczby używanych słów, czyli też $N^{O(1)}$.
  \end{itemize}
  \uncover<6->{Zatem każda liczba ma $O(\log N)$ bitów, czyli mieści się w $O(1)$ słowach modelu word RAM.}
\end{frame}

\begin{frame}{Uwaga 12: związek z innymi technikami}
  \begin{uwagaP}{12}
    \begin{itemize}
      \item Prace o mnożeniu macierzy zwykle liczą tylko \emph{mnożenia}. Tutaj dodawania trzeba kontrolować starannie.
      \item Kodowania (2.4.1) to algorytm Yatesa dla potęg Kroneckera [Yat37, Goo58], standardowy w rekurencyjnych algorytmach
        mnożenia macierzy.
      \item Przycięta rekursja (2.4.2) to analog przycinania FFT [Mar71, SB93], przyciętej inwersji Möbiusa [BHKK10] oraz
        (po transpozycji) rzadkiej ewaluacji obwodów potęg Kroneckera [AL25], opartej na [Wil24].
      \item \textbf{Nowa} w tym kierunku jest analiza z 2.4.3: przycięta rekursja odwiedza \textbf{mało liści}.
    \end{itemize}
  \end{uwagaP}
\end{frame}
